% SPDX-License-Identifier: MIT
% Compile from the repository root:
% latexmk -pdf -interaction=nonstopmode -halt-on-error -outdir=docs/build docs/user-guide.tex
\documentclass[11pt,a4paper]{article}

\usepackage[T1]{fontenc}
\usepackage[utf8]{inputenc}
\usepackage[margin=25mm]{geometry}
\usepackage{lmodern}
\usepackage{booktabs}
\usepackage{longtable}
\usepackage{array}
\usepackage{listings}
\usepackage[obeyspaces]{url}
\usepackage{xurl}
\usepackage[hidelinks]{hyperref}

\hypersetup{
  pdftitle={PIRFlow: User Guide},
  pdfauthor={João Isler},
  pdfsubject={Installation, configuration, and flow reconstruction workflows}
}
\setlength{\parindent}{0pt}
\setlength{\parskip}{0.5em}
\setlength{\emergencystretch}{2em}
\renewcommand{\arraystretch}{1.15}
\newcommand{\code}[1]{\nolinkurl{#1}}
\newcolumntype{L}[1]{>{\raggedright\arraybackslash}p{#1}}
\lstset{
  basicstyle=\ttfamily\small,
  columns=fullflexible,
  keepspaces=true,
  breaklines=true,
  showstringspaces=false,
  frame=single,
  framerule=0.3pt,
  aboveskip=0.7em,
  belowskip=0.7em
}

\title{PIRFlow\\\large User Guide}
\author{João Isler}
\date{}

\begin{document}
\maketitle

PIRFlow (Physics-Informed Reconstruction of Flow fields) reconstructs steady,
two-dimensional compressible flow fields using PyTorch. Forward runs learn
from sampled CFD fields; inverse runs learn from observations and optional
boundary data. Physics-informed runs also minimize Euler or Reynolds-averaged
Navier--Stokes (RANS) equation residuals at collocation points.

\setcounter{tocdepth}{1}
\tableofcontents

\section{Installation and quick start}
\label{sec:quick-start}

PIRFlow requires Python 3.12. On macOS or Linux, run these commands from the
repository root, where \code{pyproject.toml} is located:

\begin{lstlisting}
python3.12 -m venv .venv
source .venv/bin/activate
python -m pip install --upgrade pip
python -m pip install .
\end{lstlisting}

Copy the complete example configuration:

\begin{lstlisting}
cp configs/configuration.yaml configs/local.yaml
\end{lstlisting}

Edit \code{configs/local.yaml} for your data, physical parameters, network,
and training settings using the following sections. The supplied configuration
describes an inverse RANS case and needs your own input files; datasets are not
bundled with the project. Launch the workflow from the repository root:

\begin{lstlisting}
pirflow --config configs/local.yaml
\end{lstlisting}

The \code{--config} argument is required. Use \code{pirflow --help} to inspect
command-line options. The workflow prepares samples, splits training,
validation, and test data, builds the network, trains, evaluates held-out
data, and optionally post-processes the full CFD mesh.

\section{Input files and paths}
\label{sec:inputs}

\subsection{Directories and filenames}

The \code{paths} section contains directories; \code{files} selects the CFD file, Gmsh
geometry, and checkpoint basename. Replace these sections in your local copy
with one entry per key. For example, with commands run from the repository
root:

\begin{lstlisting}
paths:
  flow: /path/to/case/cfd
  mesh: /path/to/case/geometry
  observations: /path/to/case/observations
  results: results/my_case
  samples: samples/my_case
  model: model/my_case

files:
  flowfield: flow.vtu
  mesh: domain.geo
  model_name: reconstruction

name: my_case
\end{lstlisting}

This reads \code{/path/to/case/cfd/flow.vtu} and, when boundary sampling is
requested, \code{/path/to/case/geometry/domain.geo}. A saved checkpoint becomes
\code{model/my_case/reconstruction.pth}; \code{name} supplies the comparison-file prefix,
giving \code{results/my_case/my_case_comparison.vtk} during post-processing.

Relative paths are resolved from the directory where you launch PIRFlow,
even when the YAML file lives in \code{configs/}. Absolute paths are also accepted.
PIRFlow creates the \code{results}, \code{samples}, and \code{model} directories; input files
must already exist. Separate output directories by case to keep saved samples
and checkpoints associated with the correct geometry and settings. Keep local
paths and datasets out of commits.

\subsection{CFD fields, coordinates, and units}

Use a VTK dataset such as \code{flow.vtu} containing the computational mesh and
flow solution. The current workflow reconstructs the two velocity components
in the x-y plane. Supply dimensional coordinates in metres and dimensional
fields with these case-sensitive array names:

\begin{longtable}{@{}L{0.27\textwidth}L{0.17\textwidth}L{0.35\textwidth}L{0.11\textwidth}@{}}
\toprule
\textbf{CFD array} & \textbf{Model variable} & \textbf{Content} & \textbf{Unit} \\
\midrule
\endfirsthead
\toprule
\textbf{CFD array} & \textbf{Model variable} & \textbf{Content} & \textbf{Unit} \\
\midrule
\endhead
\midrule
\multicolumn{4}{r}{\small Continued on the next page} \\
\endfoot
\bottomrule
\endlastfoot
\code{Density} &
  \code{rho} &
  Scalar density &
  kg/m\textsuperscript{3} \\
\code{Velocity} &
  \code{u}, \code{v} &
  Vector; components 0 and 1 &
  m/s \\
\code{Pressure} &
  \code{p} &
  Scalar static pressure &
  Pa \\
\code{Eddy_Viscosity} &
  \code{mut} &
  Scalar dynamic turbulent viscosity &
  Pa\,s \\
\end{longtable}

Forward Euler sampling requires the first three arrays; forward RANS sampling
also requires \code{Eddy_Viscosity}. Inverse boundary sampling reads the fields
selected for each boundary. The CFD sampler currently uses these fixed array
names, so rename differently named arrays before sampling. Changing
\code{post_processing.fields} only changes the reference mapping for post-processing.

Full-mesh post-processing reads reference arrays from VTK point data. Provide
point-associated fields for that stage, converting cell data beforehand if
necessary. Use the same coordinate system and units for CFD, Gmsh geometry,
and CSV observations. PIRFlow applies the physical reference scales internally;
do not pre-normalize the input fields. \code{Eddy_Viscosity} must contain dynamic
turbulent viscosity, not kinematic viscosity or a viscosity ratio.

\subsection{Inputs required by each stage}

\begin{itemize}
  \item \textbf{Fresh sampling:} reads the CFD mesh to locate the domain, including for
  inverse PINN collocation points. Gradient sampling also needs the array named
  by \code{sampling.gradient_sampling.variable}.
  \item \textbf{Boundary sampling:} reads the Gmsh geometry when a boundary count is
  positive. Its physical group names must match \code{sampling.boundaries.names}
  exactly, and its geometry must align with the CFD mesh.
  \item \textbf{Inverse observations:} velocity and pressure CSVs are read from
  \code{paths.observations} only when their modalities are enabled. Synthetic
  Schlieren uses the CFD density field; its array name is set by
  \code{identification.observations.schlieren.density_name}. See the inverse
  observations section for filenames and CSV columns.
  \item \textbf{Saved samples:} \code{run.routines.sampling: false} loads the relevant NPZ
  files from \code{paths.samples}. It does not disable observation loading,
  synthetic Schlieren generation, or full-mesh post-processing; those stages
  still require their original inputs when enabled.

\end{itemize}
\section{Configuration essentials}
\label{sec:configuration}

Start from the complete \code{configs/configuration.yaml} and edit the existing
keys in your local copy. The examples below show selected sections; they are
not complete runnable configurations. Retain the other sections, since model
construction reads some settings even when the corresponding routine is off.

\subsection{Choose the workflow and equations}

\code{run.problem} chooses the source of training data; \code{run.model} controls whether
equation residuals are included. The allowed configurations are:

\begin{longtable}{@{}L{0.32\textwidth}L{0.2\textwidth}L{0.18\textwidth}L{0.2\textwidth}@{}}
\toprule
\textbf{Workflow} & \textbf{\code{run.problem}} & \textbf{\code{run.model}} & \textbf{Architecture} \\
\midrule
\endfirsthead
\toprule
\textbf{Workflow} & \textbf{\code{run.problem}} & \textbf{\code{run.model}} & \textbf{Architecture} \\
\midrule
\endhead
\midrule
\multicolumn{4}{r}{\small Continued on the next page} \\
\endfoot
\bottomrule
\endlastfoot
Fit sampled CFD fields &
  \code{forward} &
  \code{supervised} &
  \code{mlp} or \code{gnn} \\
Fit CFD with physical constraints &
  \code{forward} &
  \code{pinn} &
  \code{mlp} or \code{gnn} \\
Reconstruct from observations with physical constraints &
  \code{inverse} &
  \code{pinn} &
  \code{mlp} \\
\end{longtable}

Forward runs do not load inverse observation CSVs. Inverse runs require
nonempty training observations or enabled boundary data. Set \code{geometry.dimension}
to \code{2}. Choose \code{run.equation: euler} for \code{rho, u, v, p}, or \code{rans} to include
the turbulent-viscosity field \code{mut} and viscous/thermal terms.

\subsection{Match the network to the problem}

Set \code{network.architecture} to \code{mlp} or \code{gnn} as allowed above. An MLP takes two
coordinates and must have four outputs for Euler or five for RANS. For example,
\code{network.mlp.layers: [2, 32, 32, 32, 5]} defines three hidden layers of width 32
for a RANS case. Set its activation and dropout in \code{network.mlp}.

For GNNs, configure \code{network.gnn} (latent width, neighbors, edge attributes,
and message-passing layers); the output width is selected from the equation.
Keep \code{network.gnn.attributes.node.boundary_marker: false}, since boundary
marker features are not implemented.

Use \code{run.device: cpu} for a CPU run or \code{cuda} for an available CUDA device.
Unavailable CUDA falls back to CPU. Other device strings are passed to PyTorch
without that availability fallback and require support for the operations used.

\subsection{Sample points and split data}

Set \code{run.routines.sampling: true} on the first run. \code{sampling.data_points}
controls CFD observations; \code{sampling.collocation_points} controls locations
where PINNs evaluate equation residuals. Supervised runs skip collocation
sampling. Inverse CFD data contain only the selected boundary observations.

For example, edit the sampling entries as follows, retaining the existing
\code{gradient_sampling} settings if gradient points are requested:

\begin{lstlisting}
sampling:
  method: lhs
  boundaries:
    names: [inlet, bottom, outlet, top]
  data_points:
    interior: 1000
    gradient: 200
    boundary: [40, 100, 40, 100]
  collocation_points:
    interior: 2000
    gradient: 400
    boundary: [40, 100, 40, 100]
\end{lstlisting}

Boundary counts follow the order of \code{names}; both count lists must have that
same length. A zero count skips that boundary. Interior counts must be positive
for fresh forward-data and collocation sampling; \code{gradient: 0} disables extra
gradient points. \code{method} accepts \code{random} or \code{lhs} (Latin hypercube sampling).
Under \code{gradient_sampling}, \code{variable} selects the CFD sensor, \code{pool_factor}
sets the candidate-pool multiplier, and \code{alpha} controls preference for larger
gradients. Domain filtering can reduce the number of retained points.

Set \code{dataset.p_training_data}, \code{p_validation_data}, and \code{p_test_data} to
nonnegative percentages summing to 100, for example 80, 10, and 10. Each dataset
being split needs at least three points; validation and test sets retain at
least one point. The top-level \code{seed} controls the shuffled split and selected
sampling/noise routines; it does not configure every random generator.

\subsection{Set physical reference scales}

Use positive, case-specific values in \code{physics.reference}:

\begin{longtable}{@{}L{0.15\textwidth}L{0.65\textwidth}L{0.14\textwidth}@{}}
\toprule
\textbf{Key} & \textbf{Meaning} & \textbf{Unit} \\
\midrule
\endfirsthead
\toprule
\textbf{Key} & \textbf{Meaning} & \textbf{Unit} \\
\midrule
\endhead
\midrule
\multicolumn{3}{r}{\small Continued on the next page} \\
\endfoot
\bottomrule
\endlastfoot
\code{Lref} &
  Reference length &
  m \\
\code{U_0} &
  Reference velocity &
  m/s \\
\code{rho} &
  Reference density &
  kg/m\textsuperscript{3} \\
\code{mu} &
  Reference molecular dynamic viscosity for RANS &
  Pa\,s \\
\end{longtable}

Coordinates are divided by \code{Lref}, velocities by \code{U_0}, density by \code{rho},
pressure by \code{rho * U_0**2}, and turbulent viscosity by \code{mu}. The pressure scale
is derived from density and velocity. Keep these scales consistent when
reloading a checkpoint.

Set \code{physics.gas.gamma} for both equation choices. RANS also uses the specific
gas constant \code{physics.gas.R} in J/(kg\,K), \code{physics.turbulence.Pr} and \code{Prt},
and the Sutherland parameters \code{T0}, \code{mu0}, and \code{S} under \code{physics.sutherland}
for temperature-dependent molecular viscosity.

\subsection{Control training and post-processing}

Enable the desired optimizers under \code{optimizer.adam} and \code{optimizer.lbfgs}.
Each enabled optimizer requires positive \code{iterations} and \code{learning_rate};
Adam runs before L-BFGS when both are enabled. Adam's scheduler multiplies its
learning rate by \code{gamma} every \code{step_size} steps. For L-BFGS, \code{iterations}
counts outer steps and \code{max_iter_per_step} limits internal iterations per step.
To skip training, set both optimizers' \code{enabled} flags to \code{false}.

\code{loss.weights.data} weights CFD and boundary fields;
\code{loss.weights.observations} weights Schlieren, velocity, and pressure
measurements. \code{loss.weights.residual} weights mass (\code{f1}), x-momentum (\code{f2}),
y-momentum (\code{f3}), and energy (\code{f4}). \code{loss.print_frequency} is a positive
integer controlling how often training losses are printed.

Set \code{run.checkpoint.save_model: true} to save after training, or
\code{run.checkpoint.load_model: true} to load the configured checkpoint before
training/evaluation. Set \code{run.routines.postprocessing: true} for mesh metrics
and VTK export. In \code{post_processing}, keep \code{fields}, \code{components}, \code{latex},
\code{colormaps}, and \code{scales} aligned in the order \code{rho, p, u, v} (then \code{mut} for
RANS). For the standard CFD arrays, use:

\begin{lstlisting}
post_processing:
  fields: [Density, Pressure, Velocity, Velocity, Eddy_Viscosity]
  components: [0, 0, 0, 1, 0]
\end{lstlisting}

Retain the other post-processing settings from the template. For Euler,
remove the final turbulent-viscosity entry from every aligned list.


\section{Inverse observations}
\label{sec:observations}

Enable or disable each modality under \code{identification.observations}:

\begin{description}
  \item[Schlieren.]
    Generated from the CFD density field. Select \code{grad_x},
    \code{grad_y}, or \code{magnitude} and configure image resolution, blur,
    noise, and sampling. This loader generates synthetic observations from
    CFD data.
  \item[Velocity profiles.]
    With \code{filename: velocity_profile} and \code{n_files: 2}, provide
    these four files:
\begin{lstlisting}
velocity_profile_u_0.csv
velocity_profile_u_1.csv
velocity_profile_v_0.csv
velocity_profile_v_1.csv
\end{lstlisting}
    Each file needs headers \code{x,y,u} or \code{x,y,v} for its component.
    Indices start at zero.
  \item[Pressure taps.]
    With \code{filename: pressure_taps}, provide \code{pressure_taps.csv}
    with headers \code{x,y,p}.
\end{description}

Optional \code{identification.boundary_conditions} selects named boundaries
and observed fields (\code{rho}, \code{u}, \code{v}, \code{p}, \code{mut}) to
sample from CFD. Those boundary names must also appear in
\code{sampling.boundaries.names}. An inverse run needs nonempty training
observations or boundary data. Fresh collocation sampling still uses the CFD
mesh to define the domain.

\section{Outputs and saved models}
\label{sec:outputs}

\begin{longtable}{@{}L{0.24\textwidth}L{0.72\textwidth}@{}}
\toprule
\textbf{Location} & \textbf{Contents} \\
\midrule
\endfirsthead
\toprule
\textbf{Location} & \textbf{Contents (continued)} \\
\midrule
\endhead
\midrule
\multicolumn{2}{r}{\small Continued on the next page} \\
\endfoot
\bottomrule
\endlastfoot
\code{paths.samples} &
  \code{data_points.npz} for forward data, \code{boundary_points.npz} for
  inverse boundary data, and \code{collocation_points.npz} for PINN runs,
  as applicable. \\
\code{paths.model} &
  \code{<files.model_name>.pth} when training with
  \code{run.checkpoint.save_model: true}. \\
\code{paths.results} &
  Plots and \code{<name>_comparison.vtk} when full-mesh post-processing is
  enabled. The VTK file includes reference, predicted, and error fields. \\
Terminal &
  Training progress and evaluation metrics: MSE, RMSE, relative L2 error,
  and $R^2$. Mesh metrics use nondimensional fields; relative L2 is displayed
  as a percentage. \\
\end{longtable}

Relative L2 is reported as \code{NaN} when reference RMS is at or below the
metric threshold. $R^2$ is \code{NaN} for fewer than two values or reference
standard deviation at or below that threshold.

To evaluate a saved model without retraining, keep its configuration, input
data, and saved samples, then set:

\begin{itemize}
  \item \code{run.routines.sampling: false}
  \item \code{run.checkpoint.load_model: true}
  \item \code{optimizer.adam.enabled: false}
  \item \code{optimizer.lbfgs.enabled: false}
  \item \code{run.routines.postprocessing: true}
\end{itemize}

Use the same architecture, physical scales, sample data, and model filename.
Inverse observation inputs are loaded again. To skip PyVista field plots while
still exporting VTK, add \code{plot: false} under \code{post_processing};
these field plots otherwise run on CPU. Sampling and training plots have
separate calls.

\section{Development and tests}
\label{sec:development}

For an editable installation with test dependencies, run these commands from
the repository root with the virtual environment activated:

\begin{lstlisting}
python -m pip install -e ".[test]"
python -m pytest -v
\end{lstlisting}

To run an individual test file:

\begin{lstlisting}
python -m pytest tests/test_metrics.py -v
\end{lstlisting}

See \code{CONTRIBUTING.md} in the repository for contribution standards and
scientific validation guidance. PIRFlow is released under the MIT License;
see \code{LICENSE} for its terms.

\end{document}
